% called by main.tex
%
\chapter{Sprint 3: Data Export and System Enhancement}
\label{ch::chapter6}

\section{Introduction}

In this chapter, we will delve into the details of Sprint 3, highlighting the progress and achievements made during this phase of the project. We will present the Sprint Backlog, discuss the functionalities developed, and showcase the design and development efforts that contributed to our ongoing project success.

\section{Sprint backlog}
the table 5.1 presents the sprint 3 backlog.
\begin{table}[!htpb]
    \caption{Sprint 3 backlog.}
    \centering
    \begin{tabular}{|>{\centering\arraybackslash}m{0.05\linewidth}|>{\centering\arraybackslash}m{0.15\linewidth}|>{\centering\arraybackslash}m{0.45\linewidth}|>{\centering\arraybackslash}m{0.1\linewidth}|>{\centering\arraybackslash}m{0.15\linewidth}|} \hline 
    \textbf{ID} & \textbf{Functionality} & \textbf{User Story} & \textbf{Priority} & \textbf{Complexity} \\ \hline
         1 & Export Data & As a user, I want to export the organized data in formats like Word, CSV, or JSON so that I can use it in other applications. & High & Medium\\ \hline
         2 & Optimize Performance & As a user, I want the system to work efficiently and accurately when processing documents so that I have a good experience without delays or errors. & Medium & Medium\\ \hline
         3 & Developing the admin interface & As an admin, I want to have control over the users activities in the application by being able to see the lists of uploaded images and exported files, and if needed, I should be able to delete them. & High & High\\ \hline
    \end{tabular}
    
    \label{tab:sprint3_backlog}
\end{table}

\section{Sprint functionalities}
In this section we will discuss the different sprint 3 implemented functionalities
\subsection{Implementing export functionality }
After submitting the data, users can now choose to export it in their preferred formats. The system supports exporting in Word, CSV and JSON, providing flexibility for various use cases.
\subsubsection{System performance optimization}
Significant efforts were made to enhance system performance, particularly in the OCR scanning and text extraction processes. This optimization ensures faster and more accurate document processing.


\subsubsection{Developing the admin interface }
In this functionality, we've added an admin interface for better control over user activities. Admins can now view and manage uploaded images and exported files, including adding or removing them as needed. This ensures data integrity and organization within the system.

\subsubsection{Documentation development}

To facilitate future updates and development, comprehensive documentation was created. This documentation provides clear and detailed information, enabling other developers to understand and work on the project efficiently.

\section{Design}
in this section we will demonstrate the different dynamic and static
diagrams of the sprint 3
\newpage
\subsection{Static diagrams}
in this part we will display the use case and class diagrams of the sprint 3 .
\subsubsection{Use case diagram}
the figure 5.1 below shows the use case diagram of sprint 3 .
\begin{figure}[!htpb]
    \centering
    \includegraphics[width=17cm]{figures/Sprint 3 Use Case.png}
    \caption{Sprint 3 use case diagram.}
    \label{fig:Sprint3_Use_Case}
\end{figure}
\newpage

\subsubsection{Class diagram}
the figure 5.2 below shows class diagram of sprint 3 .
\begin{figure}[!htpb]
    \centering
    \includegraphics[width=16.7cm]{figures/Sprint 3 Class.png}
    \caption{Sprint 3 class diagram.}
    \label{fig:Sprint3_Class}
\end{figure}

\subsection{Dynamic diagrams}
This section contains the dynamic diagrams of sprint 3.

\subsection{Sequence diagram}
Since the first sprint is where we finalized our work,  we will get back in this section to the full story of interactions between the user and our system's entities from the first to the last functionality, as mapped out in the figure 5.3 below.
\begin{figure}[!htpb]
    \centering
    \includegraphics[width=16cm]{figures/Sprint 3 Sequence.drawio.png}
    \caption{Sprint 3 sequence diagram.}
    \label{fig:Sprint3_Sequence}
\end{figure}\\
Here is where we added the admin sequence diagram, it illustrates all the changes and interactions between the admin and the other entities involved. It starts by the login phase, followed by four alternative scenarios focusing respectively on images management, file management, group management and user management, ending by the logout phase. As detailed in the figure 5.4 below. 
\begin{figure}[!htpb]
    \centering
    \includegraphics[height=19cm]{figures/Global Sequence Admin.png}
    \caption{Sprint 3 sequence diagram.}
    \label{fig:Sprint3_Sequence}
\end{figure}
\newpage



\section{Development}
In this section, we will detail the development of sprint 3, where we built upon the progress made in the preceding sprints to refine and finalize our features, while also introducing new ones.

\subsection{Data exporting}
When the user clicks the "Submit" button, the ETL function is activated, storing the extracted data in three formats: Word, CSV, and JSON. Following this, the export interface appears, allowing the user to choose their desired format. By clicking the "Export" button, the files will be downloaded to the user's machine.

The figure 5.5 shows the displayed export interface that appear after clicking the "submit" button.
\begin{figure}[!htpb]
    \centering
    \includegraphics[width=13cm]{figures/Export_Phase_interface.png}
    \caption{Export interface.}
    \label{fig:Export_Phase_interface}
\end{figure}

The figure 5.6 shows the error message that appears when the user clicks "export" button without choosing any exporting format .
\begin{figure}[!htpb]
    \centering
    \includegraphics[width=13cm]{figures/export_error_interface.png}
    \caption{Export error interface.}
    \label{fig:export_error_interface}
\end{figure}

\subsection{Performance optimization}

This section discusses the performance optimization of the three main functionalities of our OCR app: Image Preprocessing, OCR Scan, and the machine learning model's API. We will outline the limitations of the initial version and present the improvements made to enhance both time and accuracy. 

The following Figure 5.7 visually represents the differences between the old and updated versions in terms of performance .
\begin{figure}[!htpb]
    \centering
    \includegraphics[height=5.7cm]{figures/Performance chart 2D New.jpg}
    \caption{Performance chart.}
    \label{fig:Performance_Chart}
\end{figure}
\newpage
\subsection{Image preprocessing}
Regarding performance optimization in Image Preprocessing, the previous version applied a static threshold function with the same values, resulting in inconsistent image quality. We improved this by implementing a dynamic function that estimates text size for each image and adjusts the threshold values accordingly. This ensures the optimal foreground-to-background ratio, enhancing image quality.

The figures 5.8 and 5.9 below compare the results of the old and optimized versions.


\begin{figure}[!htpb]
    \centering
    \begin{minipage}{0.48\textwidth}
        \centering
        \includegraphics[height=11.5cm]{figures/before_preprocessing_optimization.jpg}
        \caption{Old preprocessed image.}
        \label{fig:before_preprocessing_optimization}
    \end{minipage}
    \hfill
    \begin{minipage}{0.48\textwidth}
        \centering
        \includegraphics[height=11.5cm]{figures/Preprocessing2.jpg}
        \caption{Optimized preprocessed image.}
        \label{fig:Optimized_Preprocessed_Image}
    \end{minipage}
\end{figure}

\newpage
\subsection{OCR scan}

To improve OCR accuracy, we tested various Tesseract parameters and found that Page Segmentation Mode (PSM) 12 and OCR Engine Mode (OEM) 1 work best. PSM 12 provides the flexibility needed for diverse text layouts, making it ideal for different document types. OEM 1, the LSTM engine, is more accurate and supports modern OCR capabilities. Additionally, it allows the inclusion of languages like French and Arabic, reducing errors with non-english characters. To enhance speed, we optimized the code by eliminating loops and adopting a more efficient approach. We also removed unnecessary functions, such as autocorrect and using large datasets, which slowed down processing by an average of 10 seconds, as illustrated in the figures 5.10 and 5.11 below .

\begin{figure}[!htpb]
    \centering
    \includegraphics[width=16cm]{figures/Good OCR.png}
    \caption{OCR output after optimization.}
    \label{fig:OCR_output_after_optimization}
\end{figure}

\begin{figure}[!htpb]
    \centering
    \includegraphics[width=16cm]{figures/bad OCR.png}
    \caption{OCR output before optimization.}
    \label{fig:OCR_output_before_optimization}
\end{figure}

\newpage
\subsection{Machine learning model}
    We optimized the model's output by upgrading from "Gemini 1.0 Pro" to "Gemini 1.5 Pro." The previous version, "Gemini 1.0 Pro" was faster but manifest significant instability. For instance, it could produce a good result for a file initially, but upon refreshing, it might generate a much poorer output. The newer version, "Gemini 1.5 Pro," while slightly slower, delivers more accurate and consistent outputs, ensuring higher stability and reliability in the results.
    
    The figures 5.12 and 5.13 show the difference in output generated for the same file when refreshed using the old machine learning model, Gemini 1.0 Pro.
\begin{figure}[!htpb]
    \centering
    \includegraphics[width=12.25cm]{figures/Gemini_Old.png}
    \caption{Bad Gemini output for file "1".}
    \label{fig:Gemini_bad_output}
\end{figure}
\begin{figure}[!htpb]
    \centering
    \includegraphics[width=12.25cm]{figures/Gemini_New.png}
    \caption{Good Gemini output for file "1"}
    \label{fig:Gemini_bad_output}
\end{figure}
\newpage
\subsection{Developing the admin interface }

We developed an admin interface to introduce a new level of control and oversight, this role allows for comprehensive management of users and their activities within the system. The admin can review, add, or remove user-uploaded images and exported files. This functionality ensures data integrity.\\
The admin can also perform the same operations on the users, groups and their assigned permissions, in addition to the ability to update the data of each of these models.

%Figures 5.14, 5.15, 5.16, 5.17, 5.18, 5.19, and 5.20 highlight the main components of our admin interface :
Figure 5.14 below shows the admin login interface, where he will provide his login credentials.
\begin{figure}[!htpb]
    \centering
    \includegraphics[width=5.1cm]{figures/Admin_Login(1).png}
    \caption{Admin login interface.}
    \label{fig:Admin_Login}
\end{figure}\\
Figure 5.15 reveals the post-login interface, which contains the recent actions of the admin and the list from which he will access the management interfaces.
\begin{figure}[!htpb]
    \centering
    \includegraphics[width=16cm]{figures/post_login_admin.png}
    \caption{Admin post login interface.}
    \label{fig:Admin_Post_Login}
\end{figure}
\newpage

    Figure 5.16 captures the admin image management interface, where he will view a list containing all the images uploaded by the users, each one with its reference, preprocessed version reference and the text exported from it, the admin is also able to create, add or Delete any of them.
\begin{figure}[!htpb]
    \centering
    \includegraphics[width=17cm]{figures/Admin_Imgs.png}
    \caption{Admin uploaded images interface.}
    \label{fig:Admin_Imgs}
\end{figure}\\
Figure 5.17 and 5.18 below outline the file management interface, where he will find the list of files, each one with its path, format, export time and the unique id of its original image.
\begin{figure}[!htpb]
    \centering
    \includegraphics[width=16cm]{figures/Admin_files_1.png}
    \caption{Admin exported files interface - part 1.}
    \label{fig:Admin_files}
\end{figure}
\begin{figure}[!htpb]
    \centering
    \includegraphics[width=17cm]{figures/Admin_files_2.png}
    \caption{ Admin exported files interface - part 2.}
    \label{fig:Admin_files}
\end{figure}\\
Figure 5.19 below clearly displays the user management interface of the admin, where the admin has full control over each user and all of his properties and permissions.
\begin{figure}[!htpb]
    \centering
    \includegraphics[width=17cm]{figures/Admin_Auth_Management_users.png}
    \caption{Admin user management interface.}
    \label{fig:Admin_Auth_Management_users}
\end{figure}\\
    Figure 5.20 below highlight the admin group management interface, where he can add a group of users and control some of their common properties such as permissions.
\begin{figure}[!htpb]
    \centering
    \includegraphics[width=17cm]{figures/Admin_Auth_Management_Groups.png}
    \caption{Admin group management interface.}
    \label{fig:Admin_Auth_Management_Groups_add}
\end{figure}
\newpage
\section{Conclusion}
In this chapter, we detailed the sprint backlog, outlined key functionalities, presented the system design and discussed our development progress. These efforts have set the stage for our final sprint, where we will refine and complete the project.